\section{Ba$_5$Y$_{12}$Zn[O(SiO$_4$)]$_8$ 的结构类似物评估}

本次检索未在 OpenAlex、Crossref 或可访问的全文语料中找到 Ba$_5$Y$_{12}$Zn[O(SiO$_4$)]$_8$ 的原始晶体结构条目。检索到的“钇钡硅酸盐”主要为发光薄膜和氧基磷灰石类材料，其化学计量比、空间群和 Zn 位点占据均不同于目标化学式 \cite{gupta2021combustion,yang2021cyan}。因此，目前较为稳妥的表述是“待验证的结构类似物候选”，而非“已知同构物”。

在此前提下仍可进行分层比较。BaZn$_2$Si$_2$O$_7$ 为 ZnO$_4$--SiO$_4$ 角共享网络和 Ba 多面体提供直接参照 \cite{lin1999phase,available2020materials,kaiser2002crystal}；钇钡硅酸盐氧基磷灰石文献提供 Y\slash Ba 与孤立硅酸根共存的组成参照，但未证明二者拓扑相同 \cite{gupta2021combustion,yang2021cyan,wei2013a}。若今后获得目标相的粉末或单晶数据，应依次核对硅酸根聚合度、Zn 配位数、Y 多面体连接方式、氧空位位置和空间群，而非仅比较化学式中的元素比例。

Ba$_5$Y$_{12}$Zn[O(SiO$_4$)]$_8$ 的高 Y 含量还提示存在电荷补偿问题。形式价态计算要求额外的 O$^{2-}$ 或阳离子混合价态；任何合成尝试都须以热重分析、XPS\slash XANES 或中子衍射确认氧计量和价态。BaZn$_2$Si$_2$O$_7$ 的热致相变可为热处理窗口的设定提供参考，但不能直接外推至富 Y 体系。表~\ref{tab:analogue}\CJKecglue{}按证据层级列出比较依据。

\begin{table}[bp]
\centering
\caption{结构类似物比较的证据层级。}
\label{tab:analogue}
\begin{tabular}{P{\dimexpr 0.1835\textwidth-2\tabcolsep\relax}P{\dimexpr 0.3538\textwidth-2\tabcolsep\relax}P{\dimexpr 0.4628\textwidth-2\tabcolsep\relax}}
\toprule
    \thead{比较层级} & \thead{BaZn$_2$Si$_2$O$_7$ 的依据} & \thead{对 Ba$_5$Y$_{12}$Zn[O(SiO$_4$)]$_8$ 的含义}\\
\midrule
    硅酸根聚合度 & ZnO$_4$\slash SiO$_4$ 角共享网络 \cite{lin1999phase} & 可作为待核对的拓扑起点，不构成同构证明\\
    \addlinespace[3pt]
    阳离子配位 & Ba 五配位、Zn 四配位 \cite{available2020materials} & 须实测 Y 和 Zn 的配位，不能按离子半径推定\\
    \addlinespace[3pt]
    热稳定性 & 约 280--300~$^\circ$C 相变 \cite{thieme2015ba1} & 仅可作为热分析程序的参照\\
\bottomrule
\end{tabular}
\end{table}
